% Part: history
% Chapter: biographies
% Section: alfred-tarski

\documentclass[../../../include/open-logic-section]{subfiles}

\begin{document}

\olfileid{his}{bio}{tar}

\olsection{আলফ্রেড তার্স্কি}

\olphoto{tarski-alfred}{আলফ্রেড তার্স্কি}

আলফ্রেড তার্স্কির জন্ম ১৯০১ সালের ১৪
জানুয়ারি পোল্যান্ডের ওয়ারশতে, যা তখন
রুশ সাম্রাজ্যের অংশ ছিল। তাঁকে
“নেপোলিয়নসুলভ” বলা হতো; তিনি ছিলেন
উচ্ছল, বাগ্মী ও প্রবল উদ্যমী। বক্তৃতাতেও
তাঁর সেই উদ্যম প্রায়ই প্রকাশ পেত। একবার
বক্তৃতার সময় সিগারেট ফেলে দিতে গিয়ে
তিনি একটি বর্জ্যপাত্রে আগুন লাগিয়ে দেন;
তার পর থেকে ওই ভবনে তাঁকে আর
বক্তৃতা দিতে দেওয়া হয়নি।

ছোটবেলা থেকেই তার্স্কির ছিল জ্ঞানপিপাসা।
পরবর্তী জীবনে তিনি শিক্ষার্থীদের বলতেন,
যুক্তিবিদ্যায় একমাত্র বি গ্রেড পেয়েছিলেন
বলেই বিষয়টি পড়তে শুরু করেন। কিন্তু
তাঁর উচ্চবিদ্যালয়ের নথিতে দেখা যায়,
যুক্তিবিদ্যাসহ সব বিষয়েই তিনি এ গ্রেড
পেয়েছিলেন। ১৯১৮ থেকে ১৯২৪ সাল
পর্যন্ত ওয়ারশ বিশ্ববিদ্যালয়ে পড়েন।
তার্স্কি প্রথমে জীববিদ্যা পড়তে চেয়েছিলেন;
কিন্তু বিশ্ববিদ্যালয়টি ওয়ারশ যুক্তিবিদ্যা
ও দর্শনচর্চার কেন্দ্র হওয়ায় গণিত,
দর্শন ও যুক্তিবিদ্যায় তাঁর আগ্রহ জন্মায়।
স্তানিস্লাভ লেশ্‌নিয়েভ্‌স্কির তত্ত্বাবধানে
১৯২৪ সালে তিনি ডক্টরেট লাভ করেন।

১৯৩৯ সালে যুক্তরাষ্ট্রে যাওয়ার আগে
ওয়ারশতে মাধ্যমিক বিদ্যালয়ের শিক্ষক
হিসেবে কাজ করার সময়েই তার্স্কি তাঁর
কয়েকটি সবচেয়ে গুরুত্বপূর্ণ গবেষণা শেষ
করেন। যৌক্তিক অনুবর্তন ও যৌক্তিক সত্য
নিয়ে তাঁর রচনাগুলি এই সময়েই লেখা।
১৯৩৯ সালে বক্তৃতা সফরে তিনি যুক্তরাষ্ট্রে
ছিলেন। সেই সফরের সময়ে জার্মানি
পোল্যান্ড আক্রমণ করে; ইহুদি বংশোদ্ভূত
হওয়ায় তার্স্কি আর দেশে ফিরতে
পারেননি। তাঁর স্ত্রী ও সন্তানেরা যুদ্ধের
শেষ পর্যন্ত পোল্যান্ডে ছিলেন; পরে
তাঁরাও যুক্তরাষ্ট্রে যেতে পারেন।
তার্স্কি হার্ভার্ডে, নিউ ইয়র্কের সিটি
কলেজে এবং প্রিন্সটনের ইনস্টিটিউট
ফর অ্যাডভান্সড স্টাডিতে পড়ানোর
পরে ক্যালিফোর্নিয়া বিশ্ববিদ্যালয়ের
বার্কলি শাখায় যান। সেখানে তিনি
যুক্তিবিদ্যা ও বিজ্ঞানের পদ্ধতিবিদ্যা
নিয়ে একটি বহুবিষয়ক কর্মসূচি
প্রতিষ্ঠা করেন। ১৯৮৩ সালের ২৬
অক্টোবর ৮২ বছর বয়সে তাঁর
মৃত্যু হয়।

\begin{reading}
তার্স্কির জীবন সম্পর্কে আরও জানতে
\emph{Alfred Tarski: Life and Logic}
জীবনীটি দেখো \citep{Feferman2004}।
যৌক্তিক অনুবর্তন ও সত্য নিয়ে
তার্স্কির পথিকৃৎ রচনাগুলি ইংরেজিতে
পাওয়া যায় \citep{Tarski1983}-তে।
তাঁর সব মূল রচনা চার খণ্ডের একটি
সংকলনে প্রকাশিত হয়েছে
\citep{Tarski1981}।
\end{reading}

\end{document}
